\documentclass[UTF8,11pt,a4paper]{ctexart}
\newcommand{\DocTitle}{Git 协作手册}
\newcommand{\DocSubtitle}{克隆、编译、提交与回退}
% Shared layout for the CW2 companion documents.
\usepackage[margin=2.1cm,headheight=16pt]{geometry}
\usepackage{amsmath,amssymb,booktabs,array,tabularx,ragged2e}
\usepackage{enumitem,xcolor,fancyhdr,listings,tcolorbox}
\tcbuselibrary{breakable}
\usepackage{hyperref}
\usepackage{bookmark}

\newcommand{\DocVersion}{3.0}
\newcommand{\EngineVersion}{1.1}
\newcommand{\DocDate}{2026-11-23}
\definecolor{ink}{HTML}{1F2937}
\definecolor{accent}{HTML}{92400E}
\definecolor{muted}{HTML}{4B5563}
\definecolor{linkink}{HTML}{174678}
\definecolor{codebg}{HTML}{F5F7FA}
\hypersetup{unicode,colorlinks=true,linkcolor=linkink,urlcolor=linkink,
  filecolor=linkink,pdftitle={SOE104 CW2 \DocTitle\space v\DocVersion},
  pdfauthor={Group A--E}}
\ctexset{section={format={\Large\bfseries\color{accent}}},
  subsection={format={\large\bfseries\color{ink}}}}
\pagestyle{fancy}
\fancyhf{}
\fancyhead[L]{\small\color{muted}SOE104 CW2 · Oregon Trail}
\fancyhead[R]{\small\color{muted}\DocTitle\ v\DocVersion}
\fancyfoot[C]{\small\color{muted}\thepage}
\renewcommand{\headrulewidth}{0.4pt}
\setlength{\parskip}{0.35em}
\setlength{\emergencystretch}{1.5em}
\renewcommand{\arraystretch}{1.23}
\setlist{itemsep=3pt,topsep=3pt,parsep=0pt,leftmargin=1.7em}
\newcolumntype{Y}{>{\RaggedRight\arraybackslash}X}
\newcolumntype{L}[1]{>{\RaggedRight\arraybackslash}p{#1}}
\newtcolorbox{keybox}[1]{breakable,title={#1},colback=orange!4,
  colframe=accent,boxrule=0.55pt,arc=2pt,fonttitle=\bfseries}
\newtcolorbox{warnbox}[1]{breakable,title={#1},colback=red!3,
  colframe=red!45!black,boxrule=0.55pt,arc=2pt,fonttitle=\bfseries}
\lstset{basicstyle=\small\ttfamily,backgroundcolor=\color{codebg},
  xleftmargin=0.9em,frame=l,rulecolor=\color{linkink!45},
  breaklines=true,columns=fullflexible,keepspaces=true,showstringspaces=false,
  aboveskip=0.65em,belowskip=0.65em}
\newcommand{\code}[1]{\texttt{#1}}
\newcommand{\source}[1]{\textup{\small\textcolor{muted}{[#1]}}}
\newcommand{\docref}[2]{\href{./#1.pdf}{#2}}
\renewcommand{\maketitle}{%
  \begin{center}
    {\LARGE\bfseries CW2 \DocTitle\par}
    \vspace{0.35em}{\large\DocSubtitle\par}
    \vspace{0.75em}{\small\color{muted}Group A--E · 整理：A
      \quad\DocDate\quad 文档 v\DocVersion\par}
    \vspace{0.2em}{\small\color{muted}Engine / Marking v\EngineVersion\par}
  \end{center}
  \vspace{0.4em}}


\begin{document}
\maketitle
\thispagestyle{fancy}

\begin{keybox}{每次提交的顺序}
查看状态并同步 $\longrightarrow$ 修改 $\longrightarrow$ 编译
$\longrightarrow$ 逐文件暂存并核对差异 $\longrightarrow$ 提交和推送。
改到共享接口时先和依赖方约好合并顺序。
\end{keybox}

\section{第一次加入仓库}

从组长拿到仓库地址，确认自己有访问权限。Git Bash 和 PowerShell 都可用。
以下 \code{REPO\_URL} 和 \code{REPO\_DIR} 需换成实际地址和目录名。
\begin{lstlisting}
git --version
git clone REPO_URL
cd REPO_DIR
git remote -v
git status --short --branch
\end{lstlisting}
确认项目目录中有 \code{src/}、\code{lib/} 和 \code{docs/}。
首次提交前若 Git 提示缺身份，在当前仓库配置作者名和邮件地址。
下面两个值需换成自己的提交身份。
\begin{lstlisting}
git config user.name "YOUR_NAME"
git config user.email "YOUR_EMAIL"
\end{lstlisting}
本地项目的 JAR 已在 \code{lib/}。IDEA 配置及命令行编译见
\docref{07_ide_setup}{《07 IDEA 配置》}。

\section{开工、提交、推送}

以下以团队使用 \code{main}、处理一个小改动为例。
若团队要求分支或审核，按约定走分支。
用 \code{git status --short --branch} 确认当前分支正确、工作区无未处理改动，再同步。
\begin{lstlisting}
git pull --ff-only
\end{lstlisting}
\code{--ff-only} 遇到分叉会停下。完成改动并编译通过后，按实际文件调整下例的路径。
\begin{lstlisting}
git diff
git add -- src/TrailHunter.java
git diff --cached
git commit -m "fix: account for available food capacity"
git push
\end{lstlisting}

\clearpage
\section{推送被拒、冲突与回退}

\begin{center}\small
\begin{tabularx}{\linewidth}{@{}L{4.3cm}Y@{}}
\toprule
遇到的情况 & 怎样处理 \\
\midrule
远端领先导致推送被拒 &
先 \code{git fetch}，再看状态和本地提交。
如果本地没有独立提交，继续用 \code{git pull --ff-only}。 \\
本地与远端都有新提交 &
工作区干净且本地提交尚未发布时，和集成人确认后可用
\code{git pull --rebase}。对已共享提交优先采用组内约定的合并方式。 \\
变基冲突 &
查看 \code{git status}，与相关作者确认文件内容。
解决后逐文件 \code{git add} 再 \code{git rebase --continue}。
若要退出未完成的变基，可运行 \code{git rebase --abort}。 \\
权限或认证失败 &
核对 \code{git remote -v} 和登录账号。请求组长检查仓库权限。
认证操作通过 Git 或 GitHub 的正常提示完成。 \\
只有一处未提交改坏 &
先看 \code{git diff -- 路径}，必要时备份该文件，只修改想恢复的部分。 \\
已共享的提交需要撤销 &
与集成人核对目标提交后，用 \code{git revert} 生成新的撤销提交，
再检查编译和差异。保持共享历史可追溯。 \\
\bottomrule
\end{tabularx}
\end{center}

\section{版本记录与交付}

平台提交的源码来自 \code{src/}。保存该次已验证的提交标识和回放结果，
必要时以对应提交打 tag。平台只计最新一次，提交前核对选中的版本。
仓库中，\code{lib/} 用于编译，\code{docs/} 保存组内文档，\code{archive/} 保存有标记的历史材料。
编译输出和本机 IDE 配置由 \code{.gitignore} 排除。
暂存时还应核对是否混入凭据或无关改动。提交说明可用 \code{feat:}、\code{fix:} 或 \code{docs:}。

\clearpage
\section{看到 Git 输出后怎么判断}

\subsection{开工前看目录和分支}

在终端不知道自己到了哪一层时，先打印仓库根目录，再看当前分支和改动。
首次克隆后看到 \code{src/}、\code{lib/}、\code{docs/} 即可对照目录。

\begin{lstlisting}
git rev-parse --show-toplevel
git status --short --branch
\end{lstlisting}

\begin{center}\small
\begin{tabularx}{\linewidth}{@{}L{4.35cm}Y@{}}
\toprule
看到的输出 & 含义与下一步 \\
\midrule
\code{not a git repository} & 当前目录不在仓库。进入克隆出的目录后再查。 \\
\code{?? docs/note.tex} & 文件已经存在，Git 还没跟踪。确认为本次交付内容再暂存。 \\
\code{main...origin/main [ahead 1]} & 根据本地远端记录，当前分支多一个本地提交。
若再提交一个文档改动，直接推送可能会把两个一起带上。 \\
\bottomrule
\end{tabularx}
\end{center}

\subsection{改完后分别检查工作区、暂存区和待推送提交}

以只更新一份文档为例，按实际情况指定 LaTeX 源和重新生成的 PDF。逐条查看输出。

\begin{lstlisting}
git diff -- docs/03_plan.tex
git add -- docs/03_plan.tex docs/03_plan.pdf
git diff --cached --stat
git diff --cached --check
\end{lstlisting}

如果误把 PDF 暂存到另一项改动，执行
\code{git restore --staged -- docs/03\_plan.pdf} 只取消暂存，文件仍保留在工作区。
提交后可刷新远端记录并列出当前分支计划推送的提交。例子使用已设置上游的分支。

\begin{lstlisting}
git fetch origin
git log --oneline "@{upstream}..HEAD"
\end{lstlisting}

\noindent 清单里每一项都属于本次打算发布的范围，再执行 \code{git push}。
遇到不认识的旧提交或分叉，先核对来源和责任人。

\end{document}
